\documentclass[a4paper,11pt]{article}
\usepackage[utf8]{inputenc}
\usepackage[T1]{fontenc}
\usepackage[french]{babel}
\usepackage[left=1cm,right=1cm,top=.5cm,bottom=0cm]{geometry}
\usepackage{enumitem}
\usepackage{hyperref}
\usepackage{xcolor}
\usepackage{titlesec}
\usepackage{fontawesome5}
\usepackage{lmodern}
\usepackage[normalem]{ulem}

\definecolor{primary}{RGB}{0, 60, 120}
\hypersetup{colorlinks=true, linkcolor=primary, filecolor=primary, urlcolor=primary}
\titleformat{\section}{\large\bfseries\color{primary}\uppercase}{}{0em}{}[\titlerule]
\titlespacing{\section}{0pt}{8pt}{5pt}

\begin{document}

\begin{center}
    {\Large \textbf{Lo\"ic Aron MBASSI EWOLO}} \\ \vspace{4pt}
    \textbf{\large Ing\'enieur Logiciel Embarqu\'e -- Syst\`emes Spatiaux} \\
    \textit{\small Disponible d\`es septembre 2026 pour alternance 12 mois} \\ \vspace{4pt}
    \small
    \faMapMarker\ Cergy, France (Mobile Cannes) \quad
    \faPhone\ (+33) 7 59 52 33 98 \quad
    \faEnvelope\ \href{mailto:loicaron.mbassiewolo@ensea.fr}{loicaron.mbassiewolo@ensea.fr} \\ \vspace{2pt}
    \faLinkedin\ \href{https://www.linkedin.com/in/loic-mbassi/}{linkedin.com/in/loic-mbassi} \quad
    \faGithub\ \href{https://github.com/Nameless0l}{github.com/Nameless0l} \quad
    \faGlobe\ \href{https://mbassiloic.tech}{mbassiloic.tech}
\end{center}

\vspace{-6pt}

\noindent\small{Les syst\`emes embarqu\'es spatiaux imposent des contraintes de fiabilit\'e, de temps r\'eel et de consommation m\'emoire que seule une ma\^itrise profonde du C/C++ bas niveau et des architectures mat\'erielles permet d'atteindre. D\'evelopper et valider des logiciels de bord pour satellites -- avec les standards de s\^uret\'e de Thales Alenia Space -- est un d\'efi technique auquel ma formation ENSEA et mes projets embarqu\'es me pr\'eparent directement.}

\section{Formation}
\begin{itemize}[leftmargin=*, label=\tiny$\bullet$, nosep]
    \item \textbf{ENSEA -- \'Ecole Nationale Sup\'erieure de l'\'Electronique et de ses Applications} \hfill \textit{2025 -- 2027} \\
    \small{Double dipl\^ome d'Ing\'enieur -- \textbf{Syst\`emes Embarqu\'es}, Signal, Traitement num\'erique, IA embarqu\'ee.}
    \item \textbf{\'Ecole Polytechnique de Yaound\'e (1\textsuperscript{\`ere} \'ecole d'ing\'enieur CEMAC)} \hfill \textit{2021 -- 2025} \\
    \small{Dipl\^ome d'Ing\'enieur de Conception (Master 2) : Math\'ematiques Appliqu\'ees, \textbf{G\'enie Logiciel}, Algorithmique.}
\end{itemize}

\section{Comp\'etences Techniques}
\begin{itemize}[leftmargin=*, nosep]
    \item \small{\textbf{Embarqu\'e :} C, C++, Python, STM32, Arduino, Raspberry Pi, Nvidia Jetson, FreeRTOS, communication s\'erie}
    \item \small{\textbf{IA embarqu\'ee :} TensorFlow Lite, TinyML, optimisation de mod\`eles, inference temps r\'eel, contraintes m\'emoire}
    \item \small{\textbf{DevOps / Qualit\'e :} Docker, Git, CI/CD, tests unitaires, int\'egration continue, documentation technique}
\end{itemize}

\section{Exp\'eriences \& Projets}
\begin{itemize}[leftmargin=*, label={-}]

\item \textbf{Stage de Recherche -- INRIA Saclay (\'Equipe TRIBE, IP Paris)} \hfill \textit{Avril -- Ao\^ut 2026} \\
\textit{Plateau de Saclay, France} \hfill \textit{Python, pipeline ML, Docker, optimisation}
\vspace{-2pt}
    \begin{itemize}[leftmargin=0.4cm, nosep, label=\tiny$\bullet$]
        \item \small{Pipeline ETL et entra\^inement BERT sur corpus 50k+ documents : automatisation compl\`ete, d\'eploiement reproductible, analyse des performances et hyperparam\`etres.}
    \end{itemize}
\vspace{3pt}

\item \textbf{PROJET-TAXI -- Syst\`eme Embarqu\'e Temps R\'eel} \hfill \textit{2023 -- 2024} \\
\textit{Syst\`eme de communication s\'ecuris\'ee} \hfill \textit{STM32, C, IPC, protocoles}
\vspace{-2pt}
    \begin{itemize}[leftmargin=0.4cm, nosep, label=\tiny$\bullet$]
        \item \small{Conception et impl\'ementation d'un syst\`eme de communication inter-processus sur plateforme STM32 : protocoles temps r\'eel, gestion des interruptions, s\'erialisation des messages.}
        \item \small{S\'ecurisation des \'echanges par chiffrement AES-256-GCM et validation par tests de charge et d'int\'egrit\'e.}
    \end{itemize}
\vspace{3pt}

\item \textbf{Harmony Gloves -- Gants Intelligents IA} \hfill \textit{2024} \\
\textit{1\textsuperscript{er} prix Mountain Hub 2024} \hfill \textit{Raspberry Pi, TF Lite, capteurs, BLE}
\vspace{-2pt}
    \begin{itemize}[leftmargin=0.4cm, nosep, label=\tiny$\bullet$]
        \item \small{J'ai particip\'e \`a la conception de gants de reconnaissance gestuelle : acquisition de capteurs IMU, inf\'erence TF Lite en temps r\'eel sur Raspberry Pi, transmission BLE optimis\'ee pour la latence.}
    \end{itemize}
\vspace{3pt}

\item \textbf{Upgrade Nvidia Jetson -- Vision Embarqu\'ee} \hfill \textit{2024} \\
\textit{Laboratoire IoT ENSEA} \hfill \textit{Jetson, CUDA, OpenCV, Python}
\vspace{-2pt}
    \begin{itemize}[leftmargin=0.4cm, nosep, label=\tiny$\bullet$]
        \item \small{Migration et optimisation d'un pipeline de vision par ordinateur sur Nvidia Jetson : portage CUDA, r\'eduction de la latence d'inference de 60\%, gestion contraintes m\'emoire emb.}
    \end{itemize}

\end{itemize}

\section{Prix \& Leadership}
\begin{itemize}[leftmargin=*, nosep, label=\tiny$\bullet$]
    \item \small{\textbf{1\textsuperscript{er} prix Mountain Hub 2024} ; \textbf{1\textsuperscript{er} prix CITS National} ; \textbf{Head of Projects Cell ENSEA}}
\end{itemize}

\end{document}
